\textbf{(CAP-COLUNI 2022 -- Q19)} Considere o trecho da reportagem, “Em 2021, Amazônia Legal registra pior acumulado de desmatamento em 5 anos”, da CNN Brasil:

De janeiro até o final de julho deste ano, a Amazônia Legal acumulou a maior área desmatada neste período na história do monitoramento feito pela plataforma Terra Brasilis. A ferramenta, desenvolvida pelo Instituto Nacional de Pesquisas Espaciais (Inpe), reúne alertas e monitora o desmatamento na região da Amazônia desde 2015 – sendo que os dados do primeiro semestre começaram a ser computados apenas a partir de 2016.

Disponível em: https://www.cnnbrasil.com.br/nacional/2021/08/06/em-2021-amazonia-legal-registra-pior-acumulado-de-desmatamento-em-5-anos. Acesso em: 20 ago. 2021.

Na tabela abaixo, temos os dados, aproximados, obtidos por essa ferramenta, nos meses de janeiro a julho de 2021.

\begin{center}
\scriptsize
\begin{tabular}{c|c}
MÊS & DESMATAMENTO (Km$^2$)\\ \hline
janeiro & 86\\
fevereiro & 123\\
março & 368\\
abril & 586\\
maio & 1391\\
junho & 1062\\
julho & 1417
\end{tabular}
\end{center}

É CORRETO afirmar que o aumento percentual do desmatamento entre os meses de janeiro e julho foi de, aproximadamente (utilize duas casas decimais):

\begin{enumerate}
	\item 1878\%
	\item 1458\%
	\item 1548\%
	\item 1748\%
\end{enumerate}
